Human communication is inherently multimodal, with emotional content conveyed not only through words but also through prosodic features such as pitch, energy, and rhythm \cite{scherer2003vocal}. However, despite their remarkable text understanding capabilities \cite{brown2020language, touvron2023llama}, current large language models (LLMs) are fundamentally limited to processing text alone, rendering them unable to perceive or respond appropriately to the rich emotional information present in speech.

This limitation has significant implications for real-world applications. Consider a mental health support chatbot: a user saying ``I'm fine'' with a flat, low-energy voice conveys distress, while the same words with high pitch and energy suggest genuine contentment. Text-only LLMs cannot distinguish between these scenarios, potentially leading to inappropriate or even harmful responses. Similar challenges arise in customer service, education, and any context requiring empathetic human-AI interaction.

Recent work has explored multimodal LLMs \cite{driess2023palm,alayrac2022flamingo}, but these primarily focus on vision-language integration. Audio-augmented LLMs remain underexplored, with most prior work either using pre-trained emotion classifiers \cite{shen2023hugginggpt} or simple acoustic features \cite{moritz2023speech} without systematic investigation of what representations enable robust emotion-aware generation.

\textbf{Our Contributions.} We address this gap through a comprehensive investigation of emotion-augmented LLM training. Our work makes three key contributions:

\begin{enumerate}
    \item \textbf{Feature Representation Analysis:} Through systematic ablation studies, we demonstrate that interpretable acoustic features (pitch contour, energy dynamics, spectral characteristics) outperform state-of-the-art deep learning embeddings (WavLM \cite{chen2022wavlm}) by 40.45\% vs 27.82\% in emotion-aware text generation. We provide detailed feature importance analysis showing that spectral features contribute 25.71\% improvement while prosodic features add an additional 14.74\% through synergistic effects.

    \item \textbf{Cross-Speaker Generalization Challenge:} We identify and characterize a severe cross-speaker generalization problem in emotion-aware LLMs. Models trained on a single speaker show 38.5× performance degradation when tested on unseen speakers, even when using supposedly speaker-invariant features like relative acoustic measures. This finding has critical implications for deployment and necessitates new approaches to speaker-invariant emotion modeling.

    \item \textbf{Multi-Speaker Training Analysis:} We investigate multi-speaker training as a potential solution, demonstrating that training on 2 speakers provides statistically significant but practically insufficient improvement (16.4\% reduction in cross-speaker loss). Critically, we show that validation on seen speakers (even with different samples) fails to predict true cross-speaker performance, revealing a 10.3× gap between same-speaker validation loss and unseen-speaker test loss.
\end{enumerate}

Our analysis of 11 comprehensive experiments provides both practical guidelines for building emotion-aware LLMs and theoretical insights into the challenges of speaker-invariant emotion perception. We release our code, data pipeline, and detailed experimental logs to facilitate future research in this important direction.
